\documentclass[12pt]{article}
\usepackage[utf8]{inputenc}
\usepackage[T1]{fontenc}
\usepackage{hyperref}
\title{Northwestern Tasmanian Language Resources for K-12 Curriculum Development}
\author{Compiled by Claude Sonnet 5.5 (claude-sonnet-5-5) for Andr\'as Kornai}
\date{}
\begin{document}
\maketitle

\section{Introduction}
Northwestern Tasmanian (ISO 639-3 xpw; also called Peerapper) is an extinct Australian language of the Western Tasmanian family. It was spoken on the west coast of Tasmania, from Macquarie Harbour north to Circular Head and Robbins Island. It has no native speakers, and my searches found no school curriculum, primer or audio for it. The attestation is a handful of nineteenth-century word lists: two ``west coast'' lists by Charles Robinson and George Augustus Robinson (246 words combined), a Robbins Island list (162 words), and Allan Cunningham's Macquarie Harbour list of 1819 (222 words). A heritage curriculum would have to be built from these lists and would be mostly about the language's history, not taught in it.

\section{Primary Educational Resources}
I found no primers, textbooks or classroom materials for Northwestern Tasmanian.

\section{Linguistic and Reference Materials}
\begin{itemize}
\item Peerapper language (Wikipedia overview and list of sources): \url{https://en.wikipedia.org/wiki/Peerapper_language}
\item Western Tasmanian languages (Wikipedia): \url{https://en.wikipedia.org/wiki/Western_Tasmanian_languages}
\item N. J. B. Plomley 1976, \textit{A Word-List of the Tasmanian Aboriginal Languages} (Glottolog reference record; the book is a rare book at AIATSIS): \url{https://glottolog.org/resource/reference/id/20405}
\item Claire Bowern 2012 on the classification of the Tasmanian languages (Glottolog reference record): \url{https://glottolog.org/resource/reference/id/306230}
\item Glottolog languoid page: \url{https://glottolog.org/resource/languoid/id/nort1442}
\end{itemize}

\section{Religious and Community Resources}
My searches turned up no scripture or community-organisation resources specific to this language. The Tasmanian Aboriginal Centre's palawa kani is a reconstructed language drawing on several Tasmanian varieties. I did not confirm how much of it draws on Northwestern Tasmanian.

\section{Digital and Audio Resources}
I found no audio recordings or digital tools for this language.

\section{Curriculum Development Context}
There are only a few short word lists and no speakers. Teachers could use the sources above to teach about the language: where it was spoken, how it is known, and how the lists were collected. Any teaching of words should be done with Tasmanian Aboriginal community bodies.
\end{document}